% !PRINTED 579
\weberpage{579}

If we further set ${}^1/_2 (n - 1) = 2l$ or $= 2l + 1$, according as
$n = 4l + 1$ or $4l + 3$, we now obtain, with
exact determination of sign,

\weq{14}{\begin{aligned} A - B &= \left( \frac{k}{n} \right) \sqrt{n}, & n \equiv 1 \ (\mathrm{mod}\ 4) \\ &= i \left( \frac{k}{n} \right) \sqrt{n}, & n \equiv 3 \ (\mathrm{mod}\ 4). \end{aligned}}

\begin{center}
\S.\,172.
\end{center}

\begin{center}
\gesperrt{The periods of} ${}^1/_3 (n - 1)$ \gesperrt{and} ${}^1/_4 (n - 1)$ \gesperrt{terms.}
\end{center}

We now pass to the case $e = 3$, in which $n - 1$
must be assumed divisible by 3, so that $n = 7, 13, 19,
31, 37, 43 \ldots$ We denote by $\varrho$ an imaginary third
root of unity:

\[ \varrho = \frac{-1 + \sqrt{-3}}{2}, \]

and determine the three periods $\eta, \eta_1, \eta_2$ of ${}^1/_3 (n - 1)$
terms each, which, as we have seen in \S.\,168, are the roots
of a cyclic cubic equation. By \S.\,169 we have

\[ \psi_1(\varrho) = \overset{t}{\underset{1,\, n - 2}{\Sigma}} \varrho^{\mathrm{ind}\, t - 2\, \mathrm{ind}(t + 1)}, \]

for which, since $-2 \equiv 1 \ (\mathrm{mod}\ 3)$, we may also set

\weq{1}{\psi_1(\varrho) = \overset{t}{\underset{1,\, n - 2}{\Sigma}} \varrho^{\mathrm{ind}(t^2 + t)}}

This number $\psi_1(\varrho)$ can be represented in one of the two
forms

\weq{2}{\psi_1(\varrho) = a + b\varrho = \frac{A + b\sqrt{-3}}{2}}

in which $a, b, A$ are integers and $A = 2a - b$.

We then have

\weq{3}{\psi_1(\varrho^2) = a + b\varrho^2 = \frac{A - b\sqrt{-3}}{2},}

and the formulae (10), (15) of \S.\,169 yield

\weq{4}{\eta + \eta_1 + \eta_2 = -1}

\weq{5}{(\varrho, \eta)^3 = n \psi_1(\varrho), \quad (\varrho^2, \eta)^3 = n \psi_1(\varrho^2)}
